\chapter{Overview}
\label{ch:overview}

% Scaffold placeholder produced by /bookwright:init. The \nocite below
% exists only so the bibliography pipeline builds before real citations
% land; remove it once Part I has genuine \cite commands.
\nocite{towell2026cipher,towell2026algebraic}

This monograph develops \emph{trapdoor computing}: a paradigm for
computing on values whose meaning is hidden behind a one-way trapdoor.
A trusted machine holds the encoder and decoder; an untrusted machine
sees only total functions on opaque bit strings flowing through opaque
lookup tables. Confidentiality here is \emph{measurable} (in the sense
of quantitative information flow) rather than \emph{negligible} (in the
sense of standard cryptographic reductions), and it comes from one-way
hashing and uniform representation, not from hiding access patterns.

The book builds the cipher map abstraction from first principles
(\cref{part:abstraction}), develops its algebra and type theory
(\cref{part:algebra}), characterizes its confidentiality at both the
marginal and compositional scales (\cref{part:confidentiality}), gives
concrete constructions (\cref{part:constructions}), and closes with the
open problems that motivate the next research program
(\cref{part:frontiers}).

% This stub is replaced by the real Part I opening from
% /bookwright:design part1 and /bookwright:write.
